\documentclass[12pt,a4paper]{article}
\usepackage[top=2.5cm,bottom=2.5cm,left=2.5cm,right=2.5cm,headheight=15pt]{geometry}
\usepackage{amsmath,amssymb}
\usepackage{booktabs,array,longtable,tabularx,multirow}
\usepackage{xcolor}
\usepackage{enumitem}
\usepackage{fancyhdr}
\usepackage{titlesec}
\usepackage{mdframed}
\usepackage{tcolorbox}
\tcbuselibrary{skins,breakable}
\usepackage{hyperref}

%% ---- Colors ----
\definecolor{planblue}{RGB}{0,82,165}
\definecolor{lightblue}{RGB}{220,235,255}
\definecolor{darkorange}{RGB}{180,70,0}
\definecolor{lightorange}{RGB}{255,235,210}
\definecolor{darkgreen}{RGB}{20,110,20}
\definecolor{lightgreen}{RGB}{200,240,205}
\definecolor{darkred}{RGB}{150,20,20}
\definecolor{lightred}{RGB}{255,210,210}
\definecolor{lightgray}{RGB}{242,242,242}

%% ---- tcolorbox styles ----
\newtcolorbox{qbox}[2][]{enhanced,breakable,
  colback=lightorange,colframe=darkorange,arc=4pt,
  title={\textbf{#2}},fonttitle=\bfseries\color{white},
  attach boxed title to top left={yshift=-2mm,xshift=4mm},
  boxed title style={colback=darkorange,arc=3pt},#1}

\newtcolorbox{solbox}[2][]{enhanced,breakable,
  colback=lightblue,colframe=planblue,arc=4pt,
  title={\textbf{#2}},fonttitle=\bfseries\color{white},
  attach boxed title to top left={yshift=-2mm,xshift=4mm},
  boxed title style={colback=planblue,arc=3pt},#1}

\newtcolorbox{keybox}[2][]{enhanced,breakable,
  colback=lightgreen,colframe=darkgreen,arc=4pt,
  title={\textbf{#2}},fonttitle=\bfseries\color{white},
  attach boxed title to top left={yshift=-2mm,xshift=4mm},
  boxed title style={colback=darkgreen,arc=3pt},#1}

\newtcolorbox{warnbox}[2][]{enhanced,breakable,
  colback=lightred,colframe=darkred,arc=4pt,
  title={\textbf{#2}},fonttitle=\bfseries\color{white},
  attach boxed title to top left={yshift=-2mm,xshift=4mm},
  boxed title style={colback=darkred,arc=3pt},#1}

%% ---- Page style ----
\pagestyle{fancy}\fancyhf{}
\lhead{\color{planblue}\textbf{Reasoning with Uncertainty --- CSE 713 AI}}
\rhead{\color{darkorange}Past Papers 2020--2024}
\cfoot{\thepage}

%% ---- Section style ----
\titleformat{\section}{\Large\bfseries\color{planblue}}{}{0em}{}[\titlerule]
\titleformat{\subsection}{\large\bfseries\color{darkorange}}{}{0em}{}

\hypersetup{colorlinks=true,linkcolor=planblue,urlcolor=planblue}

%% ---- Title ----
\title{\color{planblue}\Huge\bfseries Reasoning with Uncertainty\\[0.4em]
  \Large Bayes' Theorem + Bayesian Networks\\[0.2em]
  \normalsize Past Paper Questions \& Solutions (2020--2024)\\[0.2em]
  \normalsize CSE 713 Artificial Intelligence}
\author{Study Notes --- June 2026}
\date{}

\begin{document}
\maketitle
\tableofcontents
\newpage

%% ===================================================
\section*{Exam Pattern Summary}
\addcontentsline{toc}{section}{Exam Pattern Summary}
%% ===================================================

\begin{keybox}{Key Exam Facts}
\begin{itemize}[noitemsep]
  \item \textbf{Yield:} 5/5 years, 4--9 marks per year (Group B, Q7 or Q8)
  \item \textbf{Pattern A --- Bayes' Theorem:} compute $P(E)$, posteriors $P(H_i|E)$, most likely cause
  \item \textbf{Pattern B --- BN joint distribution:} express $P(B,E,A,J,M)$, compute specific scenario
  \item \textbf{Recurring topology:} \{Fire/Burglary, Earthquake\} $\to$ Alarm $\to$ \{Caller1, Caller2\}
  \item \textbf{Key formula:} $P(H_i|E) = \dfrac{P(E|H_i)\cdot P(H_i)}{\sum_{k=1}^{K} P(E|H_k)\cdot P(H_k)}$
  \item \textbf{BN joint:} $P(X_1,\ldots,X_n) = \prod_i P(X_i|\text{Parents}(X_i))$
\end{itemize}
\end{keybox}

\begin{warnbox}{2022 WARNING --- CPT values look wrong but are correct for that paper}
2022 B-3d has P(A$|$F=T,E=T)=0.05 (alarm less likely when BOTH fire AND earthquake). This unusual CPT is verbatim from the exam paper. Use as given.\\
2023 Q7b uses P(F)=0.4, P(E)=0.6 (NOT 0.004/0.003 --- those are 2022 values).
\end{warnbox}

\newpage
%% ===================================================
\section{2024 --- Q7 (9 marks total)}
%% ===================================================

\subsection{Q7a --- Robot Vacuum Bayesian Network (4 marks)}

\begin{qbox}{2024 Q7a}
A robot vacuum cleaner suddenly stops working (Evidence $E$). There are three possible hypotheses:
\begin{itemize}[noitemsep]
  \item $H_1$: Battery is low \quad $P(H_1)=0.5$, \quad $P(E|H_1)=0.7$
  \item $H_2$: Dust bin is full \quad $P(H_2)=0.3$, \quad $P(E|H_2)=0.8$
  \item $H_3$: Motor is overheated \quad $P(H_3)=0.2$, \quad $P(E|H_3)=0.9$
\end{itemize}
\begin{enumerate}
  \item Compute $P(E)$, the total probability that the robot stops working.
  \item Calculate the posterior probabilities $P(H_1|E)$, $P(H_2|E)$, $P(H_3|E)$.
  \item Identify the most likely cause of stopping.
  \item Justify your answer.
\end{enumerate}
\end{qbox}

\begin{solbox}{Solution}
\textbf{Step 1: Compute $P(E)$ using Total Probability:}
\[
P(E) = \sum_{i=1}^{3} P(E|H_i)\cdot P(H_i)
     = (0.7)(0.5) + (0.8)(0.3) + (0.9)(0.2)
     = 0.35 + 0.24 + 0.18 = \boxed{0.77}
\]

\textbf{Step 2: Compute posteriors using Extended Bayes' Theorem:}
\[
P(H_i|E) = \frac{P(E|H_i)\cdot P(H_i)}{P(E)}
\]
\begin{align*}
P(H_1|E) &= \frac{0.7 \times 0.5}{0.77} = \frac{0.35}{0.77} \approx \boxed{0.4545}\\[4pt]
P(H_2|E) &= \frac{0.8 \times 0.3}{0.77} = \frac{0.24}{0.77} \approx \boxed{0.3117}\\[4pt]
P(H_3|E) &= \frac{0.9 \times 0.2}{0.77} = \frac{0.18}{0.77} \approx \boxed{0.2338}
\end{align*}
\textit{Sanity check:} $0.4545 + 0.3117 + 0.2338 = 1.000$ \checkmark

\textbf{Step 3 \& 4: Most likely cause and justification:}\\
$H_1$ (Battery is low) is the most likely cause with posterior probability $\approx 45.45\%$.\\
Although $H_3$ has the highest likelihood $P(E|H_3)=0.9$, the strong prior $P(H_1)=0.5$ makes $H_1$ dominate after combining both prior and likelihood via Bayes' theorem.
\end{solbox}

\subsection{Q7b --- Challenges in Bayesian Decision Theory (2 marks)}

\begin{qbox}{2024 Q7b}
Identify and discuss the challenges in Bayesian decision theory.
\end{qbox}

\begin{solbox}{Solution}
Bayesian decision theory faces several key challenges:
\begin{enumerate}
  \item \textbf{Prior elicitation:} Specifying accurate prior probability distributions is difficult. For $n$ binary variables, a complete joint distribution requires $2^n$ entries — exponential growth. Expert elicitation is subjective and error-prone.
  \item \textbf{Computational intractability:} Exact inference in general Bayesian Networks is NP-hard. Inference scales poorly with network size, necessitating approximate methods (MCMC, variational inference).
  \item \textbf{Sensitivity to priors:} Results can vary significantly with different prior specifications, especially with limited training data. This subjectivity undermines objectivity of decisions.
  \item \textbf{Conditional independence assumptions:} The BN structure encodes strong independence assumptions that may not hold in real-world domains, leading to incorrect probability estimates.
  \item \textbf{Data requirements:} Learning accurate CPTs requires large amounts of training data. With sparse data, parameter estimates are unreliable.
\end{enumerate}
\end{solbox}

\subsection{Q7c --- Evolutionary Method (3 marks)}

\begin{qbox}{2024 Q7c}
Describe the evolutionary method with an appropriate example.
\end{qbox}

\begin{solbox}{Solution}
\textbf{Evolutionary Algorithms (EAs)} are population-based metaheuristic search methods inspired by biological evolution (natural selection, genetics).

\textbf{Key Components:}
\begin{itemize}[noitemsep]
  \item \textbf{Population:} A set of candidate solutions (individuals/chromosomes)
  \item \textbf{Fitness function:} Evaluates the quality of each individual (higher = better solution)
  \item \textbf{Selection:} Choose fitter individuals for reproduction (survival of the fittest)
  \item \textbf{Crossover:} Combine two parent chromosomes to produce offspring (exploration)
  \item \textbf{Mutation:} Random small changes to maintain diversity and avoid premature convergence
\end{itemize}

\textbf{Example --- Genetic Algorithm for Robot Path Planning:}
\begin{itemize}[noitemsep]
  \item \textit{Encoding:} Chromosome = sequence of waypoints, e.g.\ $[\text{A},\text{B},\text{D},\text{G}]$
  \item \textit{Fitness:} Inverse of total path length (shorter path $\to$ higher fitness)
  \item \textit{Selection:} Tournament selection --- compare 2 random individuals, keep the better
  \item \textit{Crossover:} Single-point crossover of two path sequences to create new routes
  \item \textit{Mutation:} Randomly swap two waypoints with small probability (e.g.\ 0.01)
  \item \textit{Termination:} After $N$ generations, output chromosome with highest fitness
\end{itemize}
After many generations, the algorithm converges to a near-optimal path through the environment.
\end{solbox}

\newpage
%% ===================================================
\section{2023 --- Q7 (9 marks total)}
%% ===================================================

\subsection{Q7a --- Three Bayes' Theorem Problems (3 marks)}

\begin{qbox}{2023 Q7a}
Solve the following Bayes' Theorem problems:
\begin{enumerate}[label=(\roman*)]
  \item \textbf{Factory:} A factory has Machine A (produces 60\% of items, 3\% defect rate) and Machine B (produces 40\%, 5\% defect rate). A product is found defective. What is $P(\text{Machine A} \mid \text{defective})$?
  \item \textbf{Medical Test:} A test is 98\% accurate. Disease prevalence is 1\%. A patient tests positive. What is $P(\text{disease} \mid \text{positive})$?
  \item \textbf{Spam Filter:} Prior $P(\text{spam})=0.20$. $P(\text{``offer''} \mid \text{spam})=0.80$, $P(\text{``offer''} \mid \neg\text{spam})=0.10$. Find $P(\text{spam} \mid \text{contains ``offer''})$.
\end{enumerate}
\end{qbox}

\begin{solbox}{Solution}

\textbf{(i) Factory Problem:}\\
Let $A$ = Machine A, $B$ = Machine B, $d$ = defective.
\[
P(d) = P(d|A)\cdot P(A) + P(d|B)\cdot P(B) = (0.03)(0.6) + (0.05)(0.4) = 0.018 + 0.020 = 0.038
\]
\[
P(A \mid d) = \frac{P(d \mid A)\cdot P(A)}{P(d)} = \frac{0.018}{0.038} \approx \boxed{0.4737}
\]
Probability the defective item came from Machine A $\approx 47.37\%$.

\medskip
\textbf{(ii) Medical Test:}\\
$P(D)=0.01$, $P(+|D)=0.98$, $P(+|\neg D)=1-0.98=0.02$ (false positive rate).
\[
P(+) = (0.98)(0.01) + (0.02)(0.99) = 0.0098 + 0.0198 = 0.0296
\]
\[
P(D \mid +) = \frac{P(+|D)\cdot P(D)}{P(+)} = \frac{0.0098}{0.0296} \approx \boxed{0.3311}
\]
Only $\approx 33.1\%$ chance of actually having disease despite positive test --- base rate matters!

\medskip
\textbf{(iii) Spam Filter:}
\[
P(\text{``offer''}) = (0.80)(0.20) + (0.10)(0.80) = 0.16 + 0.08 = 0.24
\]
\[
P(\text{spam} \mid \text{``offer''}) = \frac{(0.80)(0.20)}{0.24} = \frac{0.16}{0.24} = \frac{2}{3} \approx \boxed{0.6667}
\]
\end{solbox}

\subsection{Q7b --- Alarm Bayesian Network (6 marks)}

\begin{warnbox}{CORRECTED Values for 2023 Q7b}
The 2023 paper uses P(F)=\textbf{0.4}, P(E)=\textbf{0.6}. The values 0.004/0.003 belong to the 2022 paper (Yakin/Samin network). Verified from original exam paper.
\end{warnbox}

\begin{qbox}{2023 Q7b}
A Bayesian Network models a home alarm system. Nodes: Fire (F), Earthquake (E), Alarm (A), Mr.\ X calls (X), Mr.\ Y calls (Y). The network topology is: $F \to A \leftarrow E$, $A \to X$, $A \to Y$.\\[4pt]
Given probabilities:
\begin{itemize}[noitemsep]
  \item $P(F)=0.4$, \quad $P(E)=0.6$
  \item $P(A|F,E)$: $\{TT: 0.95,\ TF: 0.90,\ FT: 0.60,\ FF: 0.01\}$
  \item $P(X|A)$: $\{T: 0.80,\ F: 0.10\}$, \quad $P(Y|A)$: $\{T: 0.70,\ F: 0.20\}$
\end{itemize}
\begin{enumerate}[label=(\roman*)]
  \item Express the full joint distribution $P(F,E,A,X,Y)$.
  \item Compute $P(A=T,\, F=F,\, E=F,\, X=T,\, Y=T)$.
  \item Compute $P(A=T,\, F=T,\, E=F,\, X=T,\, Y=T)$.
\end{enumerate}
\end{qbox}

\begin{solbox}{Solution}

\textbf{(i) Joint Distribution (Chain Rule of BN):}
\[
P(F,E,A,X,Y) = P(F)\cdot P(E)\cdot P(A \mid F,E)\cdot P(X \mid A)\cdot P(Y \mid A)
\]

\textbf{(ii) Alarm sounds, no fire, no earthquake, both X and Y call:}
\begin{align*}
&P(A{=}T,\,F{=}F,\,E{=}F,\,X{=}T,\,Y{=}T)\\
&= P(F{=}F)\cdot P(E{=}F)\cdot P(A{=}T|F{=}F,E{=}F)\cdot P(X{=}T|A{=}T)\cdot P(Y{=}T|A{=}T)\\
&= (1-0.4)\cdot(1-0.6)\cdot 0.01 \cdot 0.80 \cdot 0.70\\
&= 0.6 \times 0.4 \times 0.01 \times 0.80 \times 0.70\\
&= 0.24 \times 0.01 \times 0.56 = \boxed{0.001344}
\end{align*}

\textbf{(iii) Alarm sounds, fire occurred, no earthquake, both X and Y call:}
\begin{align*}
&P(A{=}T,\,F{=}T,\,E{=}F,\,X{=}T,\,Y{=}T)\\
&= P(F{=}T)\cdot P(E{=}F)\cdot P(A{=}T|F{=}T,E{=}F)\cdot P(X{=}T|A{=}T)\cdot P(Y{=}T|A{=}T)\\
&= 0.4 \times 0.4 \times 0.90 \times 0.80 \times 0.70\\
&= 0.16 \times 0.90 \times 0.56 = 0.16 \times 0.504 = \boxed{0.08064}
\end{align*}
\end{solbox}

\newpage
%% ===================================================
\section{2022 --- B-3d (3 marks)}
%% ===================================================

\begin{qbox}{2022 B-3d --- Alarm Network (Yakin \& Samin)}
A Bayesian Network: Fire(F), Earthquake(E), Alarm(A), Yakin calls(Y), Samin calls(S).\\
Topology: $F \to A \leftarrow E$, $A \to Y$, $A \to S$.\\[4pt]
Given:
\begin{itemize}[noitemsep]
  \item $P(F)=0.004$, \quad $P(E)=0.003$
  \item $P(A|F,E)$: $\{TT: 0.05,\ TF: 0.94,\ FT: 0.29,\ FF: 0.001\}$
  \item $P(Y|A)$: $\{T: 0.90,\ F: 0.05\}$, \quad $P(S|A)$: $\{T: 0.80,\ F: 0.01\}$
\end{itemize}
\begin{enumerate}[label=(\roman*)]
  \item Express the joint probability distribution $P(F,E,A,Y,S)$.
  \item Compute $P(A=T,\, F=F,\, E=F,\, Y=T,\, S=T)$.
  \item Compute $P(A=T,\, F=T,\, E=F,\, Y=T,\, S=T)$.
\end{enumerate}
\end{qbox}

\begin{solbox}{Solution}

\textbf{(i) Joint Distribution:}
\[
P(F,E,A,Y,S) = P(F)\cdot P(E)\cdot P(A \mid F,E)\cdot P(Y \mid A)\cdot P(S \mid A)
\]

\textbf{(ii) Alarm sounds, no fire, no earthquake, both Y and S call:}
\begin{align*}
&P(A{=}T,\,F{=}F,\,E{=}F,\,Y{=}T,\,S{=}T)\\
&= (1-0.004)\cdot(1-0.003)\cdot P(A{=}T|FF)\cdot P(Y{=}T|A{=}T)\cdot P(S{=}T|A{=}T)\\
&= 0.996 \times 0.997 \times 0.001 \times 0.90 \times 0.80\\
&= 0.993012 \times 0.001 \times 0.72 = 0.993012 \times 0.00072 \approx \boxed{7.15 \times 10^{-4}}
\end{align*}

\textbf{(iii) Alarm sounds, fire occurred, no earthquake, both Y and S call:}
\begin{align*}
&P(A{=}T,\,F{=}T,\,E{=}F,\,Y{=}T,\,S{=}T)\\
&= 0.004 \times (1-0.003) \times P(A{=}T|TF) \times P(Y{=}T|A{=}T) \times P(S{=}T|A{=}T)\\
&= 0.004 \times 0.997 \times 0.94 \times 0.90 \times 0.80\\
&= 0.003988 \times 0.94 \times 0.72 = 0.003749 \times 0.72 \approx \boxed{2.70 \times 10^{-3}}
\end{align*}
\end{solbox}

\newpage
%% ===================================================
\section{2021 --- Q7 (8.75 marks total)}
%% ===================================================

\subsection{Q7a --- Concept of Uncertainty (2.75 marks)}

\begin{qbox}{2021 Q7a}
Explain the concept of uncertainty using the doorbell-at-midnight example. The doorbell rang at 12 o'clock midnight.
\begin{itemize}[noitemsep]
  \item Was someone at the door?
  \item Did Karim wake up?
\end{itemize}
Propositions: Prop1: $\text{AtDoor}(x) \to \text{Doorbell}$; \quad Prop2: $\text{Doorbell} \to \text{Wake(Karim)}$
\end{qbox}

\begin{solbox}{Solution}
\textbf{Question 1 --- Was someone at the door?}\\
This requires abductive reasoning: since Doorbell occurred, conclude $\text{AtDoor}(x)$ via Prop1.\\
\textit{Problem:} Prop1 is \textbf{incomplete}. The doorbell might ring due to:
\begin{itemize}[noitemsep,topsep=2pt]
  \item Electrical malfunction, wind pressure, an animal, a faulty switch
  \item The list of possible causes is effectively \textbf{infinite}
\end{itemize}
Conclusion: We \textbf{cannot} conclude with certainty that someone was at the door.

\medskip
\textbf{Question 2 --- Did Karim wake up?}\\
This requires deductive reasoning: Doorbell $\to$ Wake(Karim) via Prop2.\\
\textit{Problem:} Prop2 is \textbf{not a tautology}. Karim may not wake up if:
\begin{itemize}[noitemsep,topsep=2pt]
  \item He is in deep sleep, wearing earplugs, or is away from home
\end{itemize}
Conclusion: We \textbf{cannot} conclude with certainty that Karim woke up.

\medskip
\textbf{Sources of uncertainty demonstrated:}
\begin{itemize}[noitemsep]
  \item \textit{Incomplete propositions} (Prop1 has infinite possible antecedents)
  \item \textit{Non-tautological implications} (Prop2 holds ``often'' but not always)
  \item \textit{Uncertain conclusions} from both abductive and deductive reasoning
\end{itemize}
The solution: model uncertainty explicitly using probabilities (Bayesian approach).
\end{solbox}

\subsection{Q7b --- Why Deductive \& Abductive Reasoning Are Not Sound (2.5 marks)}

\begin{qbox}{2021 Q7b}
Explain why deductive and abductive (adductive) reasoning are not sound. What is the source of uncertainty in such reasoning?
\end{qbox}

\begin{solbox}{Solution}
\textbf{Deductive reasoning is not always sound because:}\\
The KB (knowledge base) is a \textit{restricted, simplified model} of the real world. Even if the deduction is formally valid within the model, it may be wrong in reality because:
\begin{itemize}[noitemsep]
  \item The model may be incomplete or incorrect
  \item Propositions assumed universally true may have exceptions
  \item Example: ``Doorbell $\to$ Wake(Karim)'' leads to a deductively valid conclusion, but the real-world proposition has exceptions (deep sleep, absent, etc.)
\end{itemize}

\textbf{Abductive reasoning is not sound because:}\\
It commits the logical fallacy of \textbf{affirming the consequent}:
\[P \to Q,\quad \text{observed } Q \quad \therefore P\]
This is logically invalid because many different causes can produce the same observation.
\begin{itemize}[noitemsep]
  \item Example: Many things can cause the doorbell to ring, not only someone at the door
\end{itemize}

\textbf{Source of uncertainty:}
\begin{enumerate}[noitemsep]
  \item Implications may be weak (not universally true)
  \item Propositions use imprecise language (``often'', ``usually'', ``rarely'')
  \item Knowledge representation is a restricted model of reality
  \item Inference process itself introduces uncertainty (deductive results may be formally correct but real-world-wrong; abductive results are not logically guaranteed)
\end{enumerate}
\end{solbox}

\subsection{Q7c --- Meningitis: Bayes' Theorem (3.5 marks)}

\begin{qbox}{2021 Q7c}
A doctor knows that meningitis causes a patient to have a stiff neck 40\% of the time. The doctor also knows that:
\[ P(\text{meningitis}) = \frac{1}{50000}, \qquad P(\text{stiff neck}) = \frac{1}{25} \]
What is the probability that a patient with a stiff neck has meningitis?
\end{qbox}

\begin{solbox}{Solution}
Let $M$ = meningitis, $S$ = stiff neck.\\
\textbf{Given:} $P(S|M)=0.40$, $P(M)=\frac{1}{50000}$, $P(S)=\frac{1}{25}$

\textbf{Applying Bayes' Theorem:}
\[
P(M \mid S) = \frac{P(S \mid M)\cdot P(M)}{P(S)}
= \frac{0.40 \times \dfrac{1}{50000}}{\dfrac{1}{25}}
= 0.40 \times \frac{25}{50000}
= \frac{0.40}{2000}
= \boxed{\frac{1}{5000} = 0.0002}
\]

\textbf{Interpretation:} Even though stiff neck is a notable symptom of meningitis (40\% chance if you have meningitis), the disease is so rare ($1$ in $50000$) that even given the symptom, the probability of actually having meningitis is only $\mathbf{0.02\%}$. This illustrates the critical importance of \textbf{base rate (prior) in Bayesian reasoning}.
\end{solbox}

\newpage
%% ===================================================
\section{2021 --- Q8a --- CPT Count vs Joint Table (3.5 marks)}
%% ===================================================

\begin{qbox}{2021 Q8a}
To safeguard your house, you recently installed two different alarm systems by two different reputable manufacturers that use completely different sensors. Considering these facts, the Bayesian network is given below:
\[
\text{Burglary} (B) \longrightarrow \text{Alarm}_1 (A_1), \quad \text{Burglary} (B) \longrightarrow \text{Alarm}_2 (A_2)
\]
\begin{enumerate}[label=(\roman*)]
  \item How many probabilities need to be specified for its Conditional Probability Tables (CPTs)?
  \item How many probabilities would need to be given if the same joint probability distribution were specified in a joint probability table?
\end{enumerate}
\end{qbox}

\begin{solbox}{Solution}

\textbf{Network structure:} 3 nodes — Burglary (B), Alarm$_1$ (A$_1$), Alarm$_2$ (A$_2$). All are binary (True/False). B is the root; A$_1$ and A$_2$ are children of B only (independent given B, since sensors are from different manufacturers).

\medskip
\textbf{(i) CPT Method --- number of parameters:}

For each binary node, we only need to specify $P(\text{node}=T \mid \ldots)$ since $P(\text{node}=F) = 1 - P(\text{node}=T)$.

\begin{center}
\begin{tabular}{llc}
\toprule
\textbf{Node} & \textbf{Parents} & \textbf{Free parameters} \\
\midrule
$B$ & None (root) & $P(B{=}T)$ \hspace{1em} \textbf{= 1} \\
$A_1$ & $B$ & $P(A_1{=}T|B{=}T)$,\ $P(A_1{=}T|B{=}F)$ \hspace{1em} \textbf{= 2} \\
$A_2$ & $B$ & $P(A_2{=}T|B{=}T)$,\ $P(A_2{=}T|B{=}F)$ \hspace{1em} \textbf{= 2} \\
\midrule
\textbf{Total} & & $\boxed{5}$ \\
\bottomrule
\end{tabular}
\end{center}

\textbf{(ii) Full Joint Probability Table --- number of parameters:}

With 3 binary variables, the joint table has $2^3 = 8$ entries. Since all probabilities must sum to 1, we only need to specify $8 - 1 = \boxed{7}$ free parameters.

\medskip
\textbf{Conclusion:} The Bayesian Network requires only \textbf{5 parameters} instead of \textbf{7} in the joint table. This is because the BN exploits the conditional independence $A_1 \perp A_2 \mid B$ (the two alarms are independent given whether a burglary occurred). The savings grow exponentially as more nodes are added.

\begin{keybox}{General Rule}
For a BN with $n$ binary nodes each having at most $k$ binary parents:
\[
\text{CPT parameters} = \sum_i 2^{|Pa(X_i)|}
\]
\[
\text{Joint table parameters} = 2^n - 1
\]
The BN is always more compact when nodes have fewer parents than the total number of nodes.
\end{keybox}

\end{solbox}

\newpage
%% ===================================================
\section{2021 --- Q8b --- Burglary Alarm Network (5.25 marks)}
%% ===================================================

\begin{qbox}{2021 Q8b}
Consider a Bayesian Network with nodes: Burglary (B), Earthquake (E), Alarm (A), JohnCalls (J), MaryCalls (M). Topology: $B \to A \leftarrow E$, $A \to J$, $A \to M$.\\[4pt]
Given probabilities:
\begin{itemize}[noitemsep]
  \item $P(B)=0.001$, \quad $P(E)=0.002$
  \item $P(A|B,E)$: $\{TT: 0.95,\ TF: 0.94,\ FT: 0.29,\ FF: 0.001\}$
  \item $P(J|A)$: $\{T: 0.90,\ F: 0.05\}$, \quad $P(M|A)$: $\{T: 0.70,\ F: 0.01\}$
\end{itemize}
\begin{enumerate}[label=(\roman*)]
  \item Express the joint probability distribution $P(B,E,A,J,M)$.
  \item Compute $P(A=T,\, B=F,\, E=F,\, J=T,\, M=T)$.
  \item Compute $P(A=T,\, B=T,\, E=F,\, J=T,\, M=T)$.
\end{enumerate}
\end{qbox}

\begin{solbox}{Solution}

\textbf{(i) Joint Distribution:}
\[
P(B,E,A,J,M) = P(B)\cdot P(E)\cdot P(A \mid B,E)\cdot P(J \mid A)\cdot P(M \mid A)
\]

\textbf{(ii) Alarm sounds, no burglary, no earthquake, both John and Mary call:}
\begin{align*}
&P(A{=}T,\,B{=}F,\,E{=}F,\,J{=}T,\,M{=}T)\\
&= (1-0.001)\cdot(1-0.002)\cdot P(A{=}T|FF)\cdot P(J{=}T|A{=}T)\cdot P(M{=}T|A{=}T)\\
&= 0.999 \times 0.998 \times 0.001 \times 0.90 \times 0.70\\
&= 0.997002 \times 0.001 \times 0.63\\
&= 0.000997002 \times 0.63 \approx \boxed{6.28 \times 10^{-4}}
\end{align*}

\textbf{(iii) Alarm sounds, burglary occurred, no earthquake, both John and Mary call:}
\begin{align*}
&P(A{=}T,\,B{=}T,\,E{=}F,\,J{=}T,\,M{=}T)\\
&= 0.001 \times (1-0.002) \times P(A{=}T|TF) \times P(J{=}T|A{=}T) \times P(M{=}T|A{=}T)\\
&= 0.001 \times 0.998 \times 0.94 \times 0.90 \times 0.70\\
&= 0.000998 \times 0.94 \times 0.63\\
&= 0.00093812 \times 0.63 \approx \boxed{5.91 \times 10^{-4}}
\end{align*}
\end{solbox}

\newpage
%% ===================================================
\section{2020 --- Q7 (Group B)}
%% ===================================================

\subsection{Q7a --- Uncertainty Concept (Doorbell) ($\approx$1.5 marks)}

\begin{qbox}{2020 Q7a}
Explain the concept of uncertainty using an appropriate example. [Same doorbell scenario as 2021 Q7a]
\end{qbox}

\begin{solbox}{Solution}
See 2021 Q7a solution above (identical concept). The doorbell-at-midnight example illustrates:
\begin{itemize}[noitemsep]
  \item Both abductive reasoning (someone at door?) and deductive reasoning (Karim woke?) fail under uncertainty
  \item Propositions that are ``often true but not tautologies'' create uncertainty
  \item Solution: model uncertainty using probability theory (Bayesian approach)
\end{itemize}
\end{solbox}

\subsection{Q7d --- Extension of Bayes' Theorem (3 marks)}

\begin{qbox}{2020 Q7d}
Describe the extension of Bayes' theorem. Give the equations and explain the semantics of two extended theorems. Name the application areas.
\end{qbox}

\begin{solbox}{Solution}

\textbf{Extended Bayes' Theorem --- Form 1 (Multiple Hypotheses):}
\[
\boxed{P(H_i \mid E) = \frac{P(E \mid H_i)\cdot P(H_i)}{\displaystyle\sum_{k=1}^{K} P(E \mid H_k)\cdot P(H_k)}}
\]
\textbf{Semantics of Form 1:}
\begin{itemize}[noitemsep]
  \item $P(H_i|E)$: \textit{posterior} --- probability hypothesis $H_i$ is true given evidence $E$
  \item $P(E|H_i)$: \textit{likelihood} --- probability of observing $E$ if $H_i$ is true
  \item $P(H_i)$: \textit{prior} --- a priori probability of $H_i$ without specific evidence
  \item $K$: total number of mutually exclusive, exhaustive hypotheses
  \item Denominator $= P(E)$ --- total probability of observing evidence $E$ across all hypotheses
\end{itemize}

\bigskip
\textbf{Extended Bayes' Theorem --- Form 2 (Updating with New Evidence):}
\[
\boxed{P(H \mid E,\, e) = \frac{P(H \mid E)\cdot P(e \mid E,\, H)}{P(e \mid E)}}
\]
\textbf{Semantics of Form 2:}
\begin{itemize}[noitemsep]
  \item Updates belief in $H$ given existing evidence $E$ when \textit{new} evidence $e$ arrives
  \item $P(H|E)$ plays the role of the new prior (posterior from previous update)
  \item Allows \textbf{sequential/incremental Bayesian updating} --- each new piece of evidence refines the estimate
  \item The size of joint probabilities needed grows as $2^n$ for $n$ propositions
\end{itemize}

\bigskip
\textbf{Application Areas:}
\begin{enumerate}[noitemsep]
  \item \textbf{Medical diagnosis:} $P(\text{disease} \mid \text{symptoms})$ --- differential diagnosis
  \item \textbf{Spam filtering:} $P(\text{spam} \mid \text{email features})$ --- Naive Bayes classifier
  \item \textbf{Robot localization:} $P(\text{position} \mid \text{sensor readings})$ --- particle filters
  \item \textbf{Weather forecasting:} Update rain probability as new atmospheric data arrives
  \item \textbf{Fault diagnosis:} Industrial systems --- identify machine failure from sensor data
\end{enumerate}
\end{solbox}

\newpage
%% ===================================================
\section{2020 --- Q8 (Group B)}
%% ===================================================

\subsection{Q8a --- Bayesian Network Syntax and Semantics (1.5 marks)}

\begin{qbox}{2020 Q8a}
Explain the syntax and semantics of a Bayesian Network.
\end{qbox}

\begin{solbox}{Solution}

\textbf{Syntax --- Structure of a Bayesian Network:}\\
A Bayesian Network $N = (X, G, P)$ consists of:
\begin{itemize}[noitemsep]
  \item A \textbf{DAG} $G = (V, E)$: directed acyclic graph with nodes $V=\{v_1,\ldots,v_n\}$ (random variables) and directed edges $E$ (causal influences)
  \item A set of \textbf{random variables} $X$ represented by the nodes of $G$
  \item A set of \textbf{Conditional Probability Tables (CPT)}: one distribution $P(X_v \mid X_{pa(v)})$ for each variable $X_v$, conditioned on its parents $pa(v)$
  \item \textit{Root nodes} (no parents) carry \textbf{prior probability} tables; all others carry CPTs
\end{itemize}

\textbf{Semantics --- What a BN Represents:}\\
A BN encodes a compact \textbf{factorization of the full joint probability distribution}:
\[
\boxed{P(X_1, X_2, \ldots, X_n) = \prod_{i=1}^{n} P(X_i \mid \text{Parents}(X_i))}
\]
\textbf{Markov Condition:} Each node is \textit{conditionally independent} of all its non-descendants given its parents. This is the independence assumption that makes BNs tractable.

\textit{Example (5-node alarm network):}
\[
P(B,E,A,J,M) = P(B)\cdot P(E)\cdot P(A|B,E)\cdot P(J|A)\cdot P(M|A)
\]
\end{solbox}

\subsection{Q8b --- Meningitis (same as 2021 Q7c) (1.5 marks)}

\begin{qbox}{2020 Q8b}
Using Bayes' theorem: $P(M)=\frac{1}{50000}$, $P(S|M)=0.40$, $P(S)=\frac{1}{25}$. Find $P(M|S)$.
\end{qbox}

\begin{solbox}{Solution}
\[
P(M \mid S) = \frac{P(S \mid M)\cdot P(M)}{P(S)} = \frac{0.40 \times \frac{1}{50000}}{\frac{1}{25}} = 0.40 \times \frac{25}{50000} = \frac{10}{50000} = \boxed{0.0002}
\]
Only $0.02\%$ --- very low due to the extremely small prior $P(M) = 1/50000$.
\end{solbox}

\subsection{Q8c --- Complex BN: Burglary, Earthquake, Reliable \& Unreliable Friend ($\approx$6.75 marks)}

\begin{qbox}{2020 Q8c}
A Bayesian Network models a burglar alarm system. Nodes: Burglary (B), Earthquake (E), Alarm (A), ReliableFriend (R), UnreliableFriend (C).\\[4pt]
Topology: $B, E \to A$; $A \to R, C$.\\[4pt]
Given:
\begin{itemize}[noitemsep]
  \item $P(B)=0.02$, $P(E)=0.001$
  \item $P(A|\neg B,\neg E)=0.01$; $P(A|B,\neg E)=0.9$; $P(A|\neg B,E)=0.9$; $P(A|B,E)=0.99$
  \item $P(R|A)=0.002$; $P(R|\neg A)=0.999$; \quad $P(C|A)=0.8$; $P(C|\neg A)=0.1$
\end{itemize}
\begin{enumerate}[label=(\roman*)]
  \item Draw the Bayesian Network.
  \item Compute prior $P(\text{alarm is sounding}) = P(A=T)$.
  \item Compute the normalized likelihood $P(A \mid R=T, C=T)$.
  \item Compute the posterior probability.
  \item How would you modify the BN if an earthquake radio report is also available?
\end{enumerate}
\end{qbox}

\begin{solbox}{Solution}

\textbf{(i) Network Structure:}
\[
B \longrightarrow A \longleftarrow E \qquad A \longrightarrow R \qquad A \longrightarrow C
\]
(5-node DAG: B and E are roots; A has parents B and E; R and C are leaves)

\medskip
\textbf{(ii) Prior $P(A=T)$ using Total Probability:}
\begin{align*}
P(A) &= P(A|B,E)P(B)P(E) + P(A|B,\neg E)P(B)P(\neg E)\\
     &\quad + P(A|\neg B,E)P(\neg B)P(E) + P(A|\neg B,\neg E)P(\neg B)P(\neg E)\\[4pt]
     &= 0.99(0.02)(0.001) + 0.9(0.02)(0.999) + 0.9(0.98)(0.001) + 0.01(0.98)(0.999)\\
     &= 0.0000198 + 0.017982 + 0.000882 + 0.0097902\\
     &\approx \boxed{0.02867}
\end{align*}

\medskip
\textbf{(iii) Normalized likelihood $P(A|R=T, C=T)$:}\\
\textit{Note:} $P(R|A)=0.002$ is unusual (reliable friend rarely signals when alarm sounds in this model). Use as given.
\begin{align*}
P(A{=}T, R{=}T, C{=}T) &= P(A{=}T)\cdot P(R{=}T|A{=}T)\cdot P(C{=}T|A{=}T)\\
  &= 0.02867 \times 0.002 \times 0.8 = 4.587 \times 10^{-5}\\[4pt]
P(A{=}F, R{=}T, C{=}T) &= P(A{=}F)\cdot P(R{=}T|A{=}F)\cdot P(C{=}T|A{=}F)\\
  &= 0.97133 \times 0.999 \times 0.1 \approx 0.09704\\[4pt]
P(A{=}T|R{=}T,C{=}T) &= \frac{4.587\times10^{-5}}{4.587\times10^{-5} + 0.09704} \approx \boxed{0.000473}
\end{align*}

\medskip
\textbf{(iv) Posterior $P(B=T \mid A=T)$:}
\[
P(A|B) = P(A|B,E)P(E) + P(A|B,\neg E)P(\neg E) = 0.99(0.001) + 0.9(0.999) = 0.00099 + 0.8991 = 0.90009
\]
\[
P(B|A{=}T) = \frac{P(A{=}T|B)\cdot P(B)}{P(A{=}T)} = \frac{0.90009 \times 0.02}{0.02867} \approx \frac{0.018002}{0.02867} \approx \boxed{0.628}
\]

\medskip
\textbf{(v) Modified BN for Earthquake Radio Report:}\\
Add a new node \textbf{RadioReport} as an additional child of Earthquake:
\begin{itemize}[noitemsep]
  \item $E \to \text{RadioReport}$: $P(\text{Report}|E=T) \approx 0.9$; $P(\text{Report}|E=F) \approx 0.05$
\end{itemize}
The modified topology becomes: $B, E \to A \to R, C$; and also $E \to \text{RadioReport}$.\\
When RadioReport is observed, it provides additional evidence for/against Earthquake, which then propagates to update $P(A)$.
\end{solbox}

\newpage
%% ===================================================
\section*{Quick Reference Summary}
\addcontentsline{toc}{section}{Quick Reference Summary}
%% ===================================================

\begin{keybox}{Core Formulas to Memorise}
\begin{tabular}{ll}
\toprule
\textbf{Formula} & \textbf{Description}\\
\midrule
$P(H_i|E) = \dfrac{P(E|H_i)P(H_i)}{\sum_k P(E|H_k)P(H_k)}$ & Extended Bayes (multiple hypotheses)\\[8pt]
$P(H|E,e) = \dfrac{P(H|E)\cdot P(e|E,H)}{P(e|E)}$ & Sequential update form\\[8pt]
$P(x_1,\ldots,x_n) = \prod_i P(x_i|\text{Pa}(x_i))$ & BN joint factorization (chain rule)\\[4pt]
$P(E) = \sum_i P(E|H_i)\cdot P(H_i)$ & Total probability (marginalization)\\
\bottomrule
\end{tabular}
\end{keybox}

\begin{keybox}{Numeric Results to Memorise}
\begin{tabular}{ll}
\toprule
\textbf{Problem} & \textbf{Key Result}\\
\midrule
2024 Q7a: Robot vacuum P(E) & $0.77$; Most likely $H_1$ (battery) $\approx 45.45\%$\\
2023 Q7a(i): Factory defect & $P(\text{MachineA}|\text{defect}) \approx 0.4737$\\
2023 Q7a(ii): Medical test & $P(\text{disease}|+) \approx 0.3311$ (only 33\%!)\\
2023 Q7a(iii): Spam filter & $P(\text{spam}|\text{``offer''}) = 2/3 \approx 0.667$\\
2023 Q7b(i): BN scenario & $P \approx 0.001344$\\
2023 Q7b(ii): BN scenario & $P \approx 0.08064$\\
2021 Q7c / 2020 Q8b: Meningitis & $P(M|S) = 0.0002 = 0.02\%$\\
2021 Q8a: CPT count & 5 params (BN) vs 7 params (joint table)\\
2021 Q8b(ii): BN no burglary & $\approx 6.28 \times 10^{-4}$\\
2021 Q8b(iii): BN burglary & $\approx 5.91 \times 10^{-4}$\\
2022 B-3d(ii): No fire/eq & $\approx 7.15 \times 10^{-4}$\\
2022 B-3d(iii): Fire only & $\approx 2.70 \times 10^{-3}$\\
\bottomrule
\end{tabular}
\end{keybox}

\begin{keybox}{BN Inference Types (4 Patterns)}
\begin{enumerate}[noitemsep]
  \item \textbf{Diagnostic} (effects $\to$ causes): Given John calls, find $P(\text{Burglary})$
  \item \textbf{Causal} (causes $\to$ effects): Given Burglary, find $P(\text{John calls})$
  \item \textbf{Inter-causal} (between causes of common effect): Given Alarm, find $P(\text{Burglary}|\text{also Earthquake})$
  \item \textbf{Mixed}: Given John calls and no Earthquake, find $P(\text{Alarm})$
\end{enumerate}
\end{keybox}

\begin{warnbox}{Common Mistakes in Exams}
\begin{itemize}[noitemsep]
  \item \textbf{Wrong 2023 values:} 2023 Q7b uses P(F)=0.4, P(E)=0.6 (NOT 0.004/0.003). The 0.004/0.003 values are from 2022 B-3d.
  \item \textbf{Forgetting to check:} $\sum_i P(H_i|E) = 1$ after computing all posteriors
  \item \textbf{Joint distribution factorization:} Must multiply ALL factors --- prior nodes $\times$ CPT nodes $\times$ leaf nodes
  \item \textbf{Meningitis base rate:} $P(M|S) = 0.0002$, NOT $0.40$ (the $40\%$ is $P(S|M)$, not the answer)
  \item \textbf{BN chain rule:} The joint = product of $P(\text{node}|\text{parents})$ for EVERY node
\end{itemize}
\end{warnbox}

\end{document}
